\par\medskip
\noindent
11.6. Para definir el isomorfismo (11.5.1), podemos limitarnos al caso
en que $S$ es estrictamente local (11.0), de modo que $\underline M$ y
$\underline M'$ son constantes, con valores
\begin{equation}\tag{11.6.1}
M=D(T'_o),\qquad M'=D(T_o),
\end{equation}
donde $T_o$ y $T'_o$ son, respectivamente, los toros maximales de
$A_o$ y $A'_o$. Identificando asimismo $\Phi_o$ con un grupo finito
ordinario, debemos definir un isomorfismo canónico entre su componente
$\Phi_o(\ell)$ y
\begin{equation*}
\Phi(\ell)=\operatorname{coker}\bigl(
 \widetilde u_\ell:M_\ell\longrightarrow\check M'_\ell\bigr).
\end{equation*}
En realidad, definiremos directamente un isomorfismo canónico
\begin{equation}\tag{11.6.2}
\Phi_o(\ell)\simeq
 \ker\bigl(\widetilde u_\ell\otimes D_\ell:
 M_\ell\otimes D_\ell\longrightarrow
 \check M'_\ell\otimes D_\ell\bigr).
\end{equation}
Como $\Phi_o(\ell)$ es finito y $M_\ell$ y $\check M'_\ell$ tienen el
mismo rango, se deducirá primero que
$\widetilde u_\ell:M_\ell\to\check M'_\ell$ es una isogenia. Se sigue
entonces que el miembro derecho de (11.6.2) es canónicamente isomorfo a
$\Phi(\ell)$:
\begin{equation}\tag{11.6.3}
\Phi(\ell)=\operatorname{coker}\widetilde u_\ell
 \simeq
 \ker(\widetilde u_\ell\otimes D_\ell).
\end{equation}
En efecto, identificando $M_\ell$ mediante $\widetilde u_\ell$ con un
submódulo de índice finito de $\check M'_\ell$, podemos escribir el
origen y el destino de $\widetilde u_\ell\otimes D_\ell$ como
\begin{equation*}
\sisi{\Phi(\ell)=V/M_\ell,\qquad
\Phi'(\ell)=V/\check M'_\ell}%
{M_\ell\otimes D_\ell=V/M_\ell,\qquad
\check M'_\ell\otimes D_\ell=V/\check M'_\ell},
\end{equation*}
donde
\begin{equation*}
V=M_\ell\otimes\mathbb Q_\ell
 \simeq\check M'_\ell\otimes\mathbb Q_\ell.
\end{equation*}
Componiendo (11.6.2) con (11.6.3) se obtiene inmediatamente el
isomorfismo (11.5.1) que había que definir; todo se reduce, pues, a

%%%%%%%%%%  scan idx 468  |  folio 457  %%%%%%%%%%

definir el isomorfismo canónico (11.6.2).

Escribiendo $M_\ell\otimes D_\ell$ (respectivamente,
$\check M'_\ell\otimes D_\ell$) como el límite inductivo de los $M_n$
(respectivamente, de los $\check M'_n$), donde $n$ recorre las
potencias de $\ell$, equivale a definir isomorfismos canónicos
\begin{equation}\tag{11.6.4}
{}_n\Phi_o\xrightarrow{\ \sim\ }
 \ker\bigl(\widetilde u_n:M_n\longrightarrow\check M'_n\bigr),
 \qquad n=\ell^{r+1},
\end{equation}
compatibles con los morfismos de transición evidentes al variar $n$.
Aquí, por abuso de notación, $\widetilde u_n$ denota ahora el
homomorfismo obtenido reduciendo módulo $n$ el homomorfismo
$\widetilde u_\ell:M_\ell\longrightarrow\check M'_\ell$.

Para ello, consideremos el esquema en grupos ${}_nA$ sobre $S$. Por la
hipótesis de reducción semiestable, es plano y cuasifinito sobre $S$
(2.2.1). Existe una filtración canónica
\begin{equation}\tag{11.6.5}
{}_nA\supset({}_nA^o)^f\supset({}_nA^o)^t\supset0,
\end{equation}
donde el exponente $f$ (respectivamente, $t$) denota la ``parte fija''
(2.2.3) (respectivamente, la ``parte tórica'' (5.1)) del esquema en
grupos al que se aplica. En la definición de la parte tórica en 5.1
habíamos supuesto que $S$ era completo, lo cual está permitido aquí
(11.0). (Además, vimos en 6.1 que la construcción se generaliza al caso
henseliano, si se desea.) Pongamos
\begin{equation}\tag{11.6.6}
{}_n\Psi={}_nA/({}_nA^o)^f.
\end{equation}
Este es un esquema en grupos plano, separado y cuasifinito sobre $S$
(SGA 3 $\mathrm{VI}_A$ 3.2 y $\mathrm{VI}_B$ 9.2), con fibra genérica
\begin{equation*}
{}_n\Psi_K=({}_nA_K)/({}_nA_K)^f,
\end{equation*}

%%%%%%%%%%  scan idx 469  |  folio 458  %%%%%%%%%%

donde, por definición,
$({}_nA_K)^f=({}_nA^o)^f_K$ denota el $r$-ésimo componente
$(n=\ell^{r+1})$ del progrupo de Barsotti--Tate $T_\ell(A_K)^f$. La
primera de las relaciones (5.5.8) proporciona entonces un isomorfismo
canónico
\begin{equation}\tag{11.6.7}
{}_n\Psi_K\xrightarrow{\ \sim\ }\underline M_{nK},
\qquad
\underline M_n\mathrel{\mathop:}=\underline M
 \otimes_{\mathbb Z}(\mathbb Z/n\mathbb Z).
\end{equation}
De hecho, se tiene lo siguiente.

\medskip
\noindent
\underline{Lema 11.6.8.} \underline{El esquema en grupos ${}_n\Psi$
sobre $S$ definido en (11.6.6) es étale sobre $S$, y el isomorfismo
(11.6.7) se prolonga de manera única a un morfismo}
\begin{equation}\tag{11.6.8.1}
{}_n\Psi\lhook\joinrel\longrightarrow\underline M_n
\end{equation}
\underline{que es una inmersión abierta.}

Para ver que ${}_n\Psi$ es étale, en virtud de (11.6.7) solo queda
verificarlo en los puntos de su fibra especial, o, equivalentemente,
verificar que esta fibra es étale. Esto es inmediato, pues la fibra en
cuestión es
\begin{equation*}
{}_nA_o/{}_nA_o^o\subset A_o/A_o^o=\Phi_o.
\end{equation*}
Se sigue que ${}_n\Psi$ es normal y, puesto que $\underline M_n$ es
finito sobre $S$, la existencia de la prolongación (11.6.8.1) se deduce
enseguida. Es un morfismo de esquemas en grupos étale sobre $S$, de modo
que su núcleo es étale sobre $S$. Como ${}_n\Psi$, y por tanto este
núcleo, es separado, mientras que su fibra genérica es el grupo unidad
por (11.6.7), el núcleo es el grupo unidad (3.1.3). Así, (11.6.8.1) es
un monomorfismo étale y, por tanto, una inmersión abierta (EGA IV
17.9.1), lo que prueba el lema.

Pongamos ahora
\begin{equation}\tag{11.6.9}
{}_n\Phi=({}_nA)^f/({}_nA^o)^f\subset{}_n\Psi.
\end{equation}
Así, ${}_n\Phi$ es un subesquema finito abierto de ${}_n\Psi$, pues
$({}_nA)^f$

%%%%%%%%%%  scan idx 470  |  folio 459  %%%%%%%%%%

es un subesquema en grupos finito abierto de ${}_nA$. Más precisamente,
se ve inmediatamente que

\begin{equation}\tag{11.6.10}
{}_n\Phi\xrightarrow{\ \sim\ }({}_n\Psi)^f.
\end{equation}
Componiendo (11.6.9) con (11.6.8.1), obtenemos una inmersión abierta
\begin{equation}\tag{11.6.11}
{}_n\Phi\lhook\joinrel\longrightarrow\underline M_n.
\end{equation}
Por otra parte,
\begin{equation*}
{}_n\Phi\times_Ss={}_n\Psi\times_Ss={}_nA_o/{}_nA_o^o,
\end{equation*}
y, puesto que $A_o^o$ es $n$-divisible (2.2.1), el último miembro se
identifica inmediatamente con ${}_n\Phi_o$. Obtenemos, por tanto, un
isomorfismo
\begin{equation}\tag{11.6.12}
{}_n\Phi\times_Ss\xrightarrow{\ \sim\ }{}_n\Phi_o,
\end{equation}
que muestra que ${}_n\Phi$ queda caracterizado, salvo isomorfismo único,
como el esquema finito étale sobre $S$ cuya fibra especial es
${}_n\Phi_o$. Así pues, para definir el isomorfismo buscado (11.6.4)
mediante la inmersión abierta (11.6.11), basta probar lo siguiente.

\medskip
\noindent
\underline{Lema 11.6.13.} \underline{La imagen del homomorfismo
canónico} (11.6.11) \underline{es igual al núcleo de}
\begin{equation*}
\widetilde u_n:\underline M_n\longrightarrow\underline{\check M}'_n.
\end{equation*}

La compatibilidad requerida de los homomorfismos (11.6.11) con los
morfismos de transición al variar $n$ es inmediata a partir de las
definiciones. Por consiguiente, el lema 11.6.13 completará la prueba del
teorema 11.5.
